% =============================================================================
% Chapter 13: Σχεσιακή Άλγεβρα - Λυμένες Ασκήσεις
% Data Mining Study Guide - NTUA
% =============================================================================

\chapter{Σχεσιακή Άλγεβρα - Λυμένες Ασκήσεις}
\label{ch:relational-algebra-exercises}

Αυτό το κεφάλαιο περιέχει λυμένες ασκήσεις για τη σχεσιακή άλγεβρα.

\begin{infobox}[Σύνοψη Τελεστών]
\begin{center}
\begin{tabular}{cll}
\toprule
\textbf{Σύμβολο} & \textbf{Όνομα} & \textbf{Περιγραφή} \\
\midrule
$\sigma$ & Selection & Επιλογή γραμμών \\
$\pi$ & Projection & Επιλογή στηλών \\
$\cup$ & Union & Ένωση συνόλων \\
$\cap$ & Intersection & Τομή συνόλων \\
$-$ & Difference & Διαφορά συνόλων \\
$\times$ & Cartesian Product & Καρτεσιανό γινόμενο \\
$\bowtie$ & Natural Join & Φυσική σύζευξη \\
$\bowtie_\theta$ & Theta Join & Σύζευξη με συνθήκη \\
$\div$ & Division & Διαίρεση \\
\bottomrule
\end{tabular}
\end{center}
\end{infobox}

% -----------------------------------------------------------------------------
% Exercise 2.1
% -----------------------------------------------------------------------------
\section{Άσκηση 2.1: Selection και Projection}

\begin{formulabox}[Εκφώνηση]
Δίνεται ο πίνακας \texttt{Employee(emp\_id, name, department, salary)}.

Γράψτε σε σχεσιακή άλγεβρα: «Βρες τα ονόματα των υπαλλήλων του τμήματος IT με
μισθό πάνω από 50000».
\end{formulabox}

\subsection*{Λύση}

\textbf{Βήμα 1: Selection (επιλογή γραμμών)}

\textit{Γιατί:} Πρώτα φιλτράρουμε τις γραμμές που πληρούν τις συνθήκες.

$$\sigma_{department='IT' \land salary > 50000}(Employee)$$

\textbf{Βήμα 2: Projection (επιλογή στηλών)}

\textit{Γιατί:} Μετά κρατάμε μόνο τη στήλη που μας ενδιαφέρει.

$$\pi_{name}(\sigma_{department='IT' \land salary > 50000}(Employee))$$

\begin{warningbox}[Σειρά Εφαρμογής]
Η σειρά Selection $\to$ Projection είναι σημαντική για βελτιστοποίηση! Αν κάνουμε
πρώτα projection, χάνουμε τις στήλες που χρειαζόμαστε για το φιλτράρισμα.
\end{warningbox}

\textbf{Τελική Απάντηση:} $\pi_{name}(\sigma_{department='IT' \land salary > 50000}(Employee))$

% -----------------------------------------------------------------------------
% Exercise 2.2
% -----------------------------------------------------------------------------
\section{Άσκηση 2.2: Natural Join}

\begin{formulabox}[Εκφώνηση]
Δίνονται:
\begin{itemize}
    \item \texttt{Student(sid, sname, dept)}
    \item \texttt{Enrollment(sid, cid, grade)}
    \item \texttt{Course(cid, cname, credits)}
\end{itemize}

Βρείτε τα ονόματα φοιτητών και τα μαθήματα στα οποία είναι εγγεγραμμένοι.
\end{formulabox}

\subsection*{Λύση}

\textbf{Βήμα 1: Αναγνώριση κοινών γνωρισμάτων}

\begin{itemize}
    \item Student $\bowtie$ Enrollment: κοινό \texttt{sid}
    \item Enrollment $\bowtie$ Course: κοινό \texttt{cid}
\end{itemize}

\textbf{Βήμα 2: Natural Join αλυσίδα}

$$Student \bowtie Enrollment \bowtie Course$$

\textbf{Βήμα 3: Projection}

$$\pi_{sname, cname}(Student \bowtie Enrollment \bowtie Course)$$

\begin{infobox}[Πώς λειτουργεί το Natural Join]
Το $\bowtie$ συνδέει πλειάδες που έχουν \textbf{ίδιες τιμές} στα κοινά γνωρίσματα
(αυτά με το ίδιο όνομα). Οι κοινές στήλες εμφανίζονται μία φορά στο αποτέλεσμα.
\end{infobox}

% -----------------------------------------------------------------------------
% Exercise 2.3
% -----------------------------------------------------------------------------
\section{Άσκηση 2.3: Theta Join}

\begin{formulabox}[Εκφώνηση]
Δίνονται:
\begin{itemize}
    \item \texttt{Employee(eid, ename, salary, manager\_id)}
    \item \texttt{Manager(mid, mname, bonus)}
\end{itemize}

Βρείτε ζεύγη (υπάλληλος, διευθυντής) όπου ο μισθός του υπαλλήλου είναι
μεγαλύτερος από το bonus του διευθυντή του.
\end{formulabox}

\subsection*{Λύση}

\textbf{Βήμα 1: Σύνδεση υπαλλήλου με διευθυντή}

Πρώτα συνδέουμε με βάση το \texttt{manager\_id = mid}:
$$Employee \bowtie_{manager\_id = mid} Manager$$

\textbf{Βήμα 2: Φιλτράρισμα}

Εφαρμόζουμε τη συνθήκη \texttt{salary > bonus}:
$$\sigma_{salary > bonus}(Employee \bowtie_{manager\_id = mid} Manager)$$

\textbf{Βήμα 3: Projection}

$$\pi_{ename, mname}(\sigma_{salary > bonus}(Employee \bowtie_{manager\_id = mid} Manager))$$

\textbf{Εναλλακτικά (ενιαίο Theta Join):}
$$\pi_{ename, mname}(Employee \bowtie_{manager\_id = mid \land salary > bonus} Manager)$$

% -----------------------------------------------------------------------------
% Exercise 2.4
% -----------------------------------------------------------------------------
\section{Άσκηση 2.4: Set Difference}

\begin{formulabox}[Εκφώνηση]
Δίνονται:
\begin{itemize}
    \item \texttt{AllStudents(sid, sname)}
    \item \texttt{Enrollment(sid, cid)}
\end{itemize}

Βρείτε τους φοιτητές που \textbf{δεν} είναι εγγεγραμμένοι σε κανένα μάθημα.
\end{formulabox}

\subsection*{Λύση}

\textbf{Βήμα 1: Εύρεση εγγεγραμμένων φοιτητών}

$$EnrolledStudents = \pi_{sid}(Enrollment)$$

\textbf{Βήμα 2: Set Difference}

\textit{Γιατί:} Το $A - B$ δίνει τα στοιχεία που είναι στο A αλλά όχι στο B.

$$NotEnrolled = \pi_{sid}(AllStudents) - \pi_{sid}(Enrollment)$$

\textbf{Βήμα 3: Επαναφορά ονομάτων (αν χρειάζεται)}

$$\pi_{sname}(AllStudents \bowtie NotEnrolled)$$

\begin{warningbox}[Απαίτηση Συμβατότητας]
Για τελεστές συνόλων ($\cup$, $\cap$, $-$), οι σχέσεις πρέπει να έχουν:
\begin{enumerate}
    \item Ίδιο αριθμό γνωρισμάτων
    \item Συμβατά domains στις αντίστοιχες θέσεις
\end{enumerate}
\end{warningbox}

% -----------------------------------------------------------------------------
% Exercise 2.5
% -----------------------------------------------------------------------------
\section{Άσκηση 2.5: Intersection με Projection}

\begin{formulabox}[Εκφώνηση]
Δίνονται:
\begin{itemize}
    \item \texttt{CS\_Students(sid, sname)} -- Φοιτητές πληροφορικής
    \item \texttt{Math\_Students(sid, sname)} -- Φοιτητές μαθηματικών
\end{itemize}

Βρείτε τους φοιτητές που σπουδάζουν \textbf{και} πληροφορική \textbf{και} μαθηματικά.
\end{formulabox}

\subsection*{Λύση}

$$CS\_Students \cap Math\_Students$$

ή ισοδύναμα (αν θέλουμε μόνο ονόματα):

$$\pi_{sname}(CS\_Students \cap Math\_Students)$$

\begin{infobox}[Ισοδυναμία]
Η τομή μπορεί να εκφραστεί με διαφορά:
$$A \cap B = A - (A - B)$$
\end{infobox}

% -----------------------------------------------------------------------------
% Exercise 2.6
% -----------------------------------------------------------------------------
\section{Άσκηση 2.6: Division (Απλή)}

\begin{formulabox}[Εκφώνηση]
Δίνονται:
\begin{itemize}
    \item \texttt{Takes(sid, cid)} -- φοιτητής παίρνει μάθημα
    \item \texttt{Required(cid)} -- υποχρεωτικά μαθήματα: $\{C1, C2\}$
\end{itemize}

Βρείτε τους φοιτητές που έχουν πάρει \textbf{όλα} τα υποχρεωτικά μαθήματα.
\end{formulabox}

\subsection*{Λύση}

\textbf{Η Division ($\div$) απαντά σε ερωτήματα «για όλα»:}

$$Takes \div Required$$

\textbf{Βήμα-βήμα αλγόριθμος:}

\begin{enumerate}
    \item Βρες όλους τους φοιτητές: $\pi_{sid}(Takes)$
    \item Βρες τι «θα έπρεπε» να έχει πάρει κάθε φοιτητής: $\pi_{sid}(Takes) \times Required$
    \item Βρες τι «δεν έχει» πάρει: $(\pi_{sid}(Takes) \times Required) - Takes$
    \item Αφαίρεσε όσους «δεν έχουν» κάτι: $\pi_{sid}(Takes) - \pi_{sid}((\pi_{sid}(Takes) \times Required) - Takes)$
\end{enumerate}

\textbf{Παράδειγμα:}

\begin{center}
\begin{tabular}{cc}
\textbf{Takes} & \\
\toprule
sid & cid \\
\midrule
S1 & C1 \\
S1 & C2 \\
S1 & C3 \\
S2 & C1 \\
S3 & C1 \\
S3 & C2 \\
\bottomrule
\end{tabular}
\quad\quad
\begin{tabular}{c}
\textbf{Required} \\
\toprule
cid \\
\midrule
C1 \\
C2 \\
\bottomrule
\end{tabular}
\quad\quad
\begin{tabular}{c}
\textbf{Result} \\
\toprule
sid \\
\midrule
S1 \\
S3 \\
\bottomrule
\end{tabular}
\end{center}

% -----------------------------------------------------------------------------
% Exercise 2.7
% -----------------------------------------------------------------------------
\section{Άσκηση 2.7: Division (Σύνθετη)}

\begin{formulabox}[Εκφώνηση]
Δίνονται:
\begin{itemize}
    \item \texttt{Supplies(supplier, part)}
    \item \texttt{Project(pname, part)} -- έργα και τα parts που χρειάζονται
\end{itemize}

Βρείτε τους προμηθευτές που μπορούν να προμηθεύσουν \textbf{όλα} τα parts που
χρειάζεται το έργο «Rocket».
\end{formulabox}

\subsection*{Λύση}

\textbf{Βήμα 1: Βρες τα parts του Rocket}

$$RocketParts = \pi_{part}(\sigma_{pname='Rocket'}(Project))$$

\textbf{Βήμα 2: Division}

$$Supplies \div RocketParts$$

\textbf{Σε SQL (για κατανόηση):}

\begin{lstlisting}[language=SQL]
SELECT DISTINCT s1.supplier
FROM Supplies s1
WHERE NOT EXISTS (
    SELECT part FROM Project WHERE pname = 'Rocket'
    EXCEPT
    SELECT part FROM Supplies s2
    WHERE s2.supplier = s1.supplier
);
\end{lstlisting}

% -----------------------------------------------------------------------------
% Exercise 2.8
% -----------------------------------------------------------------------------
\section{Άσκηση 2.8: Σύνθετη Έκφραση}

\begin{formulabox}[Εκφώνηση]
Δίνονται:
\begin{itemize}
    \item \texttt{Employee(eid, ename, dept\_id, salary)}
    \item \texttt{Department(dept\_id, dname, budget)}
    \item \texttt{Project(pid, pname, dept\_id)}
    \item \texttt{WorksOn(eid, pid, hours)}
\end{itemize}

Βρείτε τα ονόματα υπαλλήλων του τμήματος «Research» που εργάζονται πάνω από 20 ώρες
σε κάποιο project.
\end{formulabox}

\subsection*{Λύση}

\textbf{Βήμα 1: Εύρεση dept\_id του Research}

$$ResearchDept = \sigma_{dname='Research'}(Department)$$

\textbf{Βήμα 2: Υπάλληλοι του Research}

$$ResearchEmps = Employee \bowtie ResearchDept$$

\textbf{Βήμα 3: Υπάλληλοι με >20 ώρες}

$$HighHours = \sigma_{hours > 20}(WorksOn)$$

\textbf{Βήμα 4: Συνδυασμός}

$$Result = \pi_{ename}(ResearchEmps \bowtie HighHours)$$

\textbf{Ενιαία έκφραση:}

$$\pi_{ename}(\sigma_{dname='Research'}(Department) \bowtie Employee \bowtie \sigma_{hours>20}(WorksOn))$$
